% answers_short.tex —— 2001 年·豫粤卷 · 简短答案速查
% 由 tools/make_short_answers.py 自 answers.tex 的 \gksolhead/\gkanswer 提取，
% 并与 2001/2001_数学真题_豫粤_答案要点.md 逐题比对后生成，勿手改（改源 MD 后重跑）。

\begin{gknotice}
\textbf{使用说明}：本节只给最终答案便于快速核对，完整推导见解析版。答案为编写者独立推导，\textbf{非官方标准答案}。
\end{gknotice}

\gkpart{选择题}
\begin{center}\begin{tabular}{ccccccccccccc}
\toprule
题号 & 1 & 2 & 3 & 4 & 5 & 6 & 7 & 8 & 9 & 10 & 11 & 12 \\
\midrule
答案 & C & A & D & A & B & B & B & B & C & B & D & D \\
\bottomrule
\end{tabular}\end{center}

\gkpart{填空题}
\begin{itemize}[leftmargin=2.6em]
  \item[\textbf{13.}] $4900$
  \item[\textbf{14.}] $\dfrac{16}{5}$
  \item[\textbf{15.}] $1$
  \item[\textbf{16.}] $2n(n-1)$
\end{itemize}

\gkpart{解答题}
\begin{itemize}[leftmargin=2.6em]
  \item[\textbf{17.}] $\pi$
  \item[\textbf{18.}] (1) $a=2,\ k=50$；(2) $1$
  \item[\textbf{19.}] (1) $\dfrac14$；(2) $\dfrac{\sqrt2}{2}$
  \item[\textbf{20.}] 第一问画面高为 $88\ \mathrm{cm}$、宽为 $55\ \mathrm{cm}$（此时 $\lambda=\dfrac58$）；第二问当 $\lambda\in\left[\dfrac23,\dfrac34\right]$ 时取 $\lambda=\dfrac23$。
  \item[\textbf{21.}] 见解析。
  \item[\textbf{22.}] (1) $f\left(\dfrac12\right)=\sqrt a,\ f\left(\dfrac14\right)=\sqrt[4]a$；(2) $f(x)$ 的一个周期为 $2$；(3) $0$
\end{itemize}
